\documentclass[10pt]{article}
\usepackage[margin=0.72in]{geometry}
\usepackage[T1]{fontenc}
\usepackage{lmodern}
\usepackage{microtype}
\usepackage{booktabs}
\usepackage{graphicx}
\usepackage{amsmath,amssymb}
\usepackage{xcolor}
\usepackage{hyperref}
\usepackage{enumitem}
\definecolor{llmred}{HTML}{E5482B}
\hypersetup{colorlinks=true,urlcolor=llmred,citecolor=llmred}
\setlength{\parindent}{0pt}
\setlength{\parskip}{4pt}

\title{\vspace{-1.5em}\textbf{Your LLM Capstone Title}}
\author{Student Name (Student ID) \and Student Name (Student ID)}
\date{COMP 4901B Large Language Models --- Fall 2026}

\begin{document}
\maketitle

\begin{abstract}
In no more than 250 words, state the question, method or system, tasks, controlled comparison, principal result, and conclusion. Replace all instructional text before submission.
\end{abstract}

\section{Introduction}
Define the capability, limitation, or application being studied. State the testable claim, summarize the approach, and list the main findings or contributions.

\section{Related Work}
Position the project against the closest LLM methods, systems, and evaluations. Cite reused models, prompts, datasets, tools, code, and judge models.

\section{Tasks, Data, and Models}
Describe task construction, data provenance and splits, model or API versions, access dates, licenses, and relevant limitations.

\section{Method or System Design}
Describe the intervention precisely. Include the prompting, adaptation, retrieval, agent loop, training objective, or system architecture needed to reproduce it. Identify what is reused and what the team implemented.

\section{Experimental Setup}
Specify baselines, controlled comparisons, metrics, judge setup, decoding parameters, seeds, compute, cost, and stopping rules. Explain how the evaluation could falsify the main claim.

\section{Results}
Report the main comparisons in tables or figures. Include uncertainty or repeated-run variation where appropriate. Interpret each central result.

\section{Analysis and Failure Cases}
Present qualitative examples, ablations, reliability checks, and at least one informative failure. Explain when and why the approach succeeds or fails.

\section{Limitations and Responsible Use}
Discuss external-validity limits, model or dataset bias, privacy and security concerns, dependence on changing APIs, and possible misuse.

\section{Conclusion}
Answer the research question using the evidence presented. Distinguish what the results show from what remains conjecture.

\section*{Team Contributions and Use of AI Tools}
List each member's contributions and approximate percentage. Disclose material use of generative AI or other automated tools for coding, experiment design, analysis, or writing, and describe verification.

\section*{Reproducibility and Resource Statement}
Give the main reproduction command, dependency and model versions, expected runtime, hardware, API requirements, estimated cost, and any resources that cannot be redistributed.

\bibliographystyle{plain}
\bibliography{references}

\appendix
\section{Prompts, Additional Results, and Implementation Details}
Optional material may appear here. The appendix and references do not count toward the 6--8 page main-text limit.

\end{document}
